\documentclass[a4paper]{article}
% generated by Docutils <http://docutils.sourceforge.net/>
\usepackage{fixltx2e} % LaTeX patches, \textsubscript
\usepackage{cmap} % fix search and cut-and-paste in Acrobat
\usepackage{ifthen}
\usepackage[T1]{fontenc}
\usepackage[utf8]{inputenc}
\usepackage{amsmath}
\usepackage{color}
\usepackage{float} % float configuration
\floatplacement{figure}{H} % place figures here definitely
\usepackage{graphicx}
\setcounter{secnumdepth}{0}
\usepackage{longtable,ltcaption,array}
\setlength{\extrarowheight}{2pt}
\newlength{\DUtablewidth} % internal use in tables

%%% Custom LaTeX preamble
% PDF Standard Fonts
\usepackage{mathptmx} % Times
\usepackage[scaled=.90]{helvet}
\usepackage{courier}

%%% User specified packages and stylesheets

%%% Fallback definitions for Docutils-specific commands

% admonition (specially marked topic)
\providecommand{\DUadmonition}[2][class-arg]{%
  % try \DUadmonition#1{#2}:
  \ifcsname DUadmonition#1\endcsname%
    \csname DUadmonition#1\endcsname{#2}%
  \else
    \begin{center}
      \fbox{\parbox{0.9\textwidth}{#2}}
    \end{center}
  \fi
}
% numeric or symbol footnotes with hyperlinks
\providecommand*{\DUfootnotemark}[3]{%
  \raisebox{1em}{\hypertarget{#1}{}}%
  \hyperlink{#2}{\textsuperscript{#3}}%
}
\providecommand{\DUfootnotetext}[4]{%
  \begingroup%
  \renewcommand{\thefootnote}{%
    \protect\raisebox{1em}{\protect\hypertarget{#1}{}}%
    \protect\hyperlink{#2}{#3}}%
  \footnotetext{#4}%
  \endgroup%
}

% title for topics, admonitions, unsupported section levels, and sidebar
\providecommand*{\DUtitle}[2][class-arg]{%
  % call \DUtitle#1{#2} if it exists:
  \ifcsname DUtitle#1\endcsname%
    \csname DUtitle#1\endcsname{#2}%
  \else
    \smallskip\noindent\textbf{#2}\smallskip%
  \fi
}

% hyperlinks:
\ifthenelse{\isundefined{\hypersetup}}{
  \usepackage[colorlinks=true,linkcolor=blue,urlcolor=blue]{hyperref}
  \urlstyle{same} % normal text font (alternatives: tt, rm, sf)
}{}
\hypersetup{
  pdftitle={Creating Good Graphics and Archiving Your Data},
}

%%% Title Data
\title{\phantomsection%
  Creating Good Graphics and Archiving Your Data%
  \label{creating-good-graphics-and-archiving-your-data}}
\author{}
\date{}

%%% Body
\begin{document}
\maketitle


\section{Introduction%
  \label{introduction}%
}

You've spent days/months/years collecting/analyzing your research topic.
You put the requested information into your paper/presentation, turn it
in, and viola, you're done, right? Not really. What you've just
generated is an electronic document of which the only remaining part may
be a hard copy in a few years. Further, without all of the splendid data
that you used to generate your charts they are now most likely
uncorrectable.

This is the case for a large portion of the theses and project reports
that students and professionals develop. Many times the only result of a
significant amount of effort is a digital photograph of a specific
representation of the data. Changing the representation of the data in
the future isn't a simple task unless a few simple steps were taken to
assure their permanency. Further, setting your graphics in the wrong
format significantly degrades the presentation of the effort on paper,
on screen, or both. This short document illustrates simple steps that
must be taken in order to archive your work for future access.


\section{Data storage%
  \label{data-storage}%
}

Plotting your data and saving the plot is not archiving the data. It is
converting it to a visual representation, one that you chose. It may not
be the one that you, or the people you send it to, want in the future.
Archiving means saving the raw data. Every amount of data processing
that moves away from raw data destroys information that is almost always
irreversible. Calculated quantities such as mean, variance, integrals,
etc, do not retain all of the information in the raw data. Further,
saving only those quantities raises an error in their calculation into a
monumental blunder should their calculation have somehow been faulty. It
is important to archive your data in the least processed form possible.
If the processing is very time consuming, archiving the processed result
is also a good idea. However, rarely is the cost of archiving raw data
excessive.


\section{How to store your data%
  \label{how-to-store-your-data}%
}

Data should always be stored in raw ASCII (American Standard Code for
Information Interchange), or raw text, files. Other formats can
\emph{additionally} be used. However, those formats may become unreadable in
10-20 years. The ASCII format has been around for a long time, and is
the format of the programming languages used to write the codes that
read other formats, so it's a good bet that this format will exist
indefinitely. A good compromise are the \href{http://docutils.sourceforge.net/rst.html}{reST} and \href{http://daringfireball.net/projects/markdown/}{Markdown}. The advantages of
these are they are ASCII, but very readable. I have more experience
with \href{http://docutils.sourceforge.net/rst.html}{reST} and favor it's ability to
be re-rendered in a slew of other formats using tools such as \href{http://sphinx-doc.org}{Sphinx}. In fact, the source of \href{http://cecs.wright.edu/~jslater/_sources/archivingdata/archivingdata.txt}{this very page}. Further,
while perusing my site you may notice that there is a link to the
left that almost always allows you to see the original source in \href{http://docutils.sourceforge.net/rst.html}{reST}
format.

Data should be stored in column delimited format. You can use whatever
character you like for delimiting. Spaces, commas, tabs, all can work
fine, but make sure that none of those characters are part of your data.
Almost all computer codes can save to an ASCII file of some type. This
file will almost always be larger than when using other formats.
Further, your program may complain that you will lose information in
saving in this format. Often, this is useless information. However, try
to find out what it is. In many cases you should also save in native
format for the code. In fact, you should almost always save in \emph{both}
formats. The ASCII format is your emergency backup for when you can no
longer read a data file written by a code that has either a) been
extinct for many years, or b) is unavailable to you due to licenses at
your new place of employment.\DUfootnotemark{id3}{id18}{1} If there is a \emph{standard} format for
your field, you should also save in that format. This means up to three
archived data files in the place of one. If this seems like a lot, then
you don't think much of your data, and this document is irrelevant.


\section{Documenting your data%
  \label{documenting-your-data}%
}

Just having your data stored in files is not enough. It must be
documented. This documentation may be in your report/thesis. If so, this
document must be part of the data archive. Further, this documentation
must be provided in ASCII format too! We really don't know if MS Word
2000 will be readable in the year 2020. Probably not. If that s the
case, when you come back to this data in the future to see if some new
processing technique yields new and exciting insight, you'd better hope
that the data is documented in an ASCII file. Typically, this file is
named \textquotedbl{}README.TXT\textquotedbl{}. The reason should be readily apparent. This is the
first document you should look at. Sets of data described by this
document should be in the same directory as this document, or
subdirectories beneath it. This document must contain:
\newcounter{listcnt0}
\begin{list}{\arabic{listcnt0}.}
{
\usecounter{listcnt0}
\setlength{\rightmargin}{\leftmargin}
}

\item A description of the experiment, including equipment used (down to
serial numbers!), names of files corresponding to pictures useful for
recreating the experiment, and enough information to find final
published documents using this data (the actual citation is
preferred). If the data is the output of a code, a version number and
the input information should be supplied. If at all possible, the
code itself should be supplied, as well as the compiler requirements
(e.g.~version 4.0 of Matlab, or LaTeX, or MS Excel 95)

\item Nomenclature for directories/subdirectories and files so that the
circumstances (case information) of each file can be ascertained.

\item Your name and where you think you are going. Some way to contact you.

\item Names of codes that can read this data (for each file type). This can
also be complex. As much as possible, presume the reader has no clue
what to do to use this data. Most likely, you will be the reader in
the future.
\end{list}


\section{Archiving your results: File Formats%
  \label{archiving-your-results-file-formats}%
}

The results of your collected data are the computational analysis and
resulting plots and figures resulting from them. No doubt this took a
great deal of effort. If all of it was generated manually, then you have
made life hard on yourself. Most tasks involving data processing
plotting are easily automated with modern tools. In performing many
tasks by hand, you require the future user of the data to know each and
every step you followed to process then plot the data. Guess what, this
is the second part of the \emph{README.TXT} file. Most likely, a user will want
to change one or two simple steps of the process. The more complex the
process, the more specific the instructions will be.


\section{Plotting%
  \label{plotting}%
}

Languages such as Octave, Gnuplot, Matlab, Mathematica, IDL, etc allow
scripting. A compiled language can also be used with the same results.
You should provide a file that when executed will generate the files
that became your plots. If the data processing is extensive, the
processing and plotting (or rendering) should be separated into two or
more tasks, with two or more scripts. Ideally these scripts should not
prompt the user for information, but define answers to pertinent
questions as variables that can be edited. Better yet is to have one
script for each data file, with the names easily identifiable. For
example: \emph{plotfrfdata.m} and \emph{frfdata.txt}.

These files need not read directly from the ASCII file, but should make
clear how to load data from the text file if necessary.


\section{File formats for graphics%
  \label{file-formats-for-graphics}%
}

Use of the wrong graphics formats for data presentation is the best way
to make good results look sloppy. The following is an incomplete list of
graphics file formats. It lists the most reliable ones for the tasks
stated. Reliable means commonly acceptable and readable.


\section{Photographs%
  \label{photographs}%
}

TIFF (Tagged Image File Format) is a very, very high quality format for
graphics such as scanned images and photographs. It also results in some
very large files. Joint Photographic Experts Group (JPEG or JPG) allows
significant compression with minimal loss of the same types of data,
however the compression can lose some information and thus jpeg
shouldn’t be used when image processing may later be performed. Both
work well when various shades of color/brightness are used, with
relatively soft edges. Storing a photograph in any other format likely
makes no sense at all. Using these formats for graphics other than
scanned images or photographs is not recommended. A very coarse
description of these formats are that for each pixel they record the
appropriate color and darkness to plot. JPG uses schemes to compress
this data. In the compression process, JPG (before 2000, which is far
form standard) uses techniques which cause strange aberrations in data
with lines or text.
\setcounter{listcnt0}{0}
\begin{list}{\arabic{listcnt0}.}
{
\usecounter{listcnt0}
\setlength{\rightmargin}{\leftmargin}
}

\item Preferred archive format; TIFF, JPG (high quality)

\item Preferred presentation format: JPG
\end{list}


\section{Line art, plots, graphs, computer renderings of objects with computer drawn axes/lines, text%
  \label{line-art-plots-graphs-computer-renderings-of-objects-with-computer-drawn-axes-lines-text}%
}

Vector graphic formats are ideal for these cases. Encapsulated
PostScript (EPS) is the best way to go, with Portable Document Format
(PDF) a close second. Windows Meta File (WMF) may work if you are
confident that you are getting vector graphics. Encapsulated PostScript
(EPS) or Portable Document Format (PDF) done correctly doesn't know what
a pixel is. Thus it prints to infinite resolution by telling the output
device simply that there is a line between two locations (not pixels)
and that the output device should adjust pixels in its resolution to
make it as perfect as possible. This differs from JPG/TIFF which can
alias and look quite bad as they are viewed closer or blown up in size.
EPS and PDF can be scaled up or down without loss in quality. Generating
a JPG or Graphics Interchange Format (GIF) and converting it to EPS is
\emph{not} equivalent and not recommended. All of the vector information is
destroyed in generating the JPG/TIFF/PNG(Portable Network Graphic)/GIF
and cannot be recovered. Screen captures currently also result in poor
quality graphics. There are some exceptions. If you think you have the
exception, print your graphic at 100 times the capture size. If the text
looks jagged to your eye at any level on a 300 dpi printer, you don't
have a vector graphic.

For example, observe Figures %
\raisebox{1em}{\hypertarget{id5}{}}\hyperlink{id4}{\textbf{\color{red}:numref:`\%s <figure1>`}} and %
\raisebox{1em}{\hypertarget{id7}{}}\hyperlink{id6}{\textbf{\color{red}:numref:`\%s <figure2>`}}. They are
sections of complex graphics. Figure %
\raisebox{1em}{\hypertarget{id9}{}}\hyperlink{id8}{\textbf{\color{red}:numref:`\%s <figure1>`}} was saved as a JPG,
then imported and edited in an unknown package. Figure %
\raisebox{1em}{\hypertarget{id11}{}}\hyperlink{id10}{\textbf{\color{red}:numref:`\%s <figure2>`}} is
the one I regenerated. Note the vastly improved cleanliness of the line.
By saving as a vector graphic, the line in  Figure %
\raisebox{1em}{\hypertarget{id13}{}}\hyperlink{id12}{\textbf{\color{red}:numref:`\%s <figure2>`}} is now
editable by hand or using an editor such as Adobe Illustrator. Figure
%
\raisebox{1em}{\hypertarget{id15}{}}\hyperlink{id14}{\textbf{\color{red}:numref:`\%s <figure1>`}} cannot easily be edited to improve the
appearance of of the line.

\DUadmonition[system-message]{
\DUtitle[system-message]{system-message}
\raisebox{1em}{\hypertarget{id4}{}}

{\color{red}ERROR/3} in \texttt{archivingdata.rst}, line~204

\hyperlink{id5}{
Unknown interpreted text role \textquotedbl{}numref\textquotedbl{}.
}}

\DUadmonition[system-message]{
\DUtitle[system-message]{system-message}
\raisebox{1em}{\hypertarget{id6}{}}

{\color{red}ERROR/3} in \texttt{archivingdata.rst}, line~204

\hyperlink{id7}{
Unknown interpreted text role \textquotedbl{}numref\textquotedbl{}.
}}

\DUadmonition[system-message]{
\DUtitle[system-message]{system-message}
\raisebox{1em}{\hypertarget{id8}{}}

{\color{red}ERROR/3} in \texttt{archivingdata.rst}, line~204

\hyperlink{id9}{
Unknown interpreted text role \textquotedbl{}numref\textquotedbl{}.
}}

\DUadmonition[system-message]{
\DUtitle[system-message]{system-message}
\raisebox{1em}{\hypertarget{id10}{}}

{\color{red}ERROR/3} in \texttt{archivingdata.rst}, line~204

\hyperlink{id11}{
Unknown interpreted text role \textquotedbl{}numref\textquotedbl{}.
}}

\DUadmonition[system-message]{
\DUtitle[system-message]{system-message}
\raisebox{1em}{\hypertarget{id12}{}}

{\color{red}ERROR/3} in \texttt{archivingdata.rst}, line~204

\hyperlink{id13}{
Unknown interpreted text role \textquotedbl{}numref\textquotedbl{}.
}}

\DUadmonition[system-message]{
\DUtitle[system-message]{system-message}
\raisebox{1em}{\hypertarget{id14}{}}

{\color{red}ERROR/3} in \texttt{archivingdata.rst}, line~204

\hyperlink{id15}{
Unknown interpreted text role \textquotedbl{}numref\textquotedbl{}.
}}

\setlength{\DUtablewidth}{\linewidth}
\begin{longtable}[c]{|p{0.179\DUtablewidth}|p{0.342\DUtablewidth}|p{0.179\DUtablewidth}|}
\caption{How to generate vector graphics.}\\
\hline
\textbf{%
Language
} & \textbf{%
Command
} & \textbf{%
Help
} \\
\hline
\endfirsthead
\caption[]{How to generate vector graphics. (... continued)}\\
\hline
\textbf{%
Language
} & \textbf{%
Command
} & \textbf{%
Help
} \\
\hline
\endhead
\multicolumn{3}{c}{\hfill ... continued on next page} \\
\endfoot
\endlastfoot

Matlab
 & 
print -depsc2 filename
 & 
(help print)
 \\
\hline

Mathematica
 & 
Export{[}\textquotedbl{}filename.eps\textquotedbl{},\textquotedbl{}eps\textquotedbl{}{]}
 & 
Help menu
 \\
\hline

MS Windows
 & 
Install a PS driver, print
with this driver
 & 
Google it
sorry. Edit
text file to
start with \%!
 \\
\hline

Macintosh
 & 
Print, note options at
bottom left of dialog box
 & 
PDF is
built-in
 \\
\hline
\end{longtable}
\label{table}

\begin{figure}
\phantomsection\label{figure1}
\noindent\makebox[\textwidth][c]{\includegraphics[width=0.500\linewidth]{image.jpg}}
\caption{Zoomed line from a JPG file. Zoom in a factor of 10 times and
notice the jaggedness of the line.}
\end{figure}

\begin{figure}
\phantomsection\label{figure2}
\noindent\makebox[\textwidth][c]{\includegraphics[width=0.500\linewidth]{image.pdf}}
\caption{Zoomed line from an EPS file. Zoom in as much as you want. The line will be no more jagged than the resolution of your screen.}
\end{figure}

The website \url{http://vectormagic.stanford.edu/} can convert pixelated
graphics to vector graphics, but there’s no substitute for the real
thing.


\section{Drawings%
  \label{drawings}%
}

The PNG and GIF formats are ideal for graphics with swaths of uniform
colors. PNG is preferred by some because there are no license
restrictions on the format, but GIF is more common. However, the
licenses and concerns that some had for the GIF format ended with the
expiration of all patents on the technology in August, 2006. Either of
these do well, resulting in high quality graphics with small file sizes.
The problem with them is that lines and text can also alias. In most
cases it is best to simply use EPS in these cases as well. EPS stores
the drawing as a set of objects that can be edited as an object instead
of as a painting. PNG and GIF do not understand objects so much as
pixels.


\section{Combined%
  \label{combined}%
}

When you have text on a photograph, what do you do? Use EPS again (or
PDF). The JPG or TIFF is embedded in the file as an object, and the text
is embedded on top of it. With these formats you can combine the best of
both worlds.


\section{Archiving Data: Storage%
  \label{archiving-data-storage}%
}

Archiving is more than how the data is organized. It also included
physically where the data (the actual storage media) is put and under
what conditions it is stored. Choice of media is of critical importance.
All of the preceding advice is eventually worthless if the media on
which you store your information fails. In my observation most computer
users who are prepared for media failure have suffered from failure
before. Those who haven't often have too much faith that their hard
drive/tapes/floppies/CDs will last forever.


\section{Archiving Media%
  \label{archiving-media}%
}

Magnetic-based media such as floppy disks and hard drives have a half
life of between five and seven years depending on storage conditions.
Technically, that means the strength of the magnetic fields drop by 50\%
in that time. Practically, it means there is a significant chance of
data loss in that time. Optical media tends to be more robust, but
varies depending on multiple factors.

Floppy disks tend to fair the worst because they are typically
manufactured as a commodity item where quality is of minimal concern.
This wasn't always the case, but it certainly is now. Floppy disk drives
are also manufactured cheaply and their read/write quality varies. Any
data currently on floppy disks (zip disks included) should be expected
to last no more than a couple of years reliably (if the data is
critical, this is insufficient). My advice is to either throw out the
floppy disks, or copy the information to a better medium, then throw out
the disks. My experience with old floppy disks is mixed.

Hard drives tend to fair better, being mostly reliable for the four or
five year life of the typical computer. However, storing something on a
hard drive on your computer is not archiving. You can archive to an
external hard drive, but it is critical that the hard drive not be
accessible again except in an emergency. It cannot remain accessible
electronically to any computer. If it does, it is vulnerable to viruses,
power surges, user errors, etc. Hard drives do not provide sufficient
long term reliability given that their cost exceeds more reliable
methods.

Tapes tend to be the most reliable of magnetic media. Life expectancies
of 10-30 years are common, but only under ideal conditions. The National
Bureau of Standards publication, Care and Handling of Computer Magnetic
Storage Media, recommends that magnetic tape be stored at 65 +/- 3
degrees Fahrenheit and 40\% +/- 5\% Relative Humidity. For whatever
reason, I’ve observed that tapes are rarely used by individuals for
archiving. This was perhaps erroneous until the advent of optical media.

Optical media can be much more reliable, but it is hard to produce, and
the cheaper media is not more reliable than magnetic media. CDs/DVDs can
fade over time. Life of the media depends on the sealing method,
reflective layer, organic dye used, and user storage conditions.
Re-writable media should no be used for archiving as they are
specifically designed to be changeable after the first burn, i.e.~the
“burn is not permanent”. Thus they are also more vulnerable to decay
over time and more susceptible to damage. Further, different formats,
e.g.~CD-R, CD+R, DVD-R and DVD+R use unique algorithms/concepts for
tracking data on the media and providing redundancy. DVD+R is the most
robust format, coming after all of the others, and being designed to
overcome the shortcomings of the others. McFarland has a rather long and
detailed document on this, and recommends using only Taiyo Yuden DVD+R
at the current time, available from the SuperMediaStore.com, purported
to be the only online US distributor that guarantees that their media is
certified Taiyo Yuden. The Taiyo Yuden FAQ also contains useful
information. Although this citation is a blog, given the depth of
discussion and responses, I believe it warrants respect. Unfortunately I
haven't looked for alternative citations on this topic. However, my
recollection of articles during the development of DVD writable media
agree with the substantial part of this document.


\section{Usage of Graphics in Documents%
  \label{usage-of-graphics-in-documents}%
}

Shortly after reading this document, you will find out that MS
PowerPoint~will not accept EPS files for show. They may print. I don't
know. However, since that’s not the intended use of MS PowerPoint, it
doesn't matter. My answer to that is that this document covers
\emph{archiving}, not usage. Although there is overlap, and the information
and explanations within still apply, practical limitations always take
precedence. For the sake of usage in MS PowerPoint, I convert my EPS
files to PNG or GIF format. In the future, JPG2000 format may be a
better choice, but compatibility is a major issue when generating a
presentation.


\section{Care of Media, Transcription, and Redundancy%
  \label{care-of-media-transcription-and-redundancy}%
}

Media has a finite storage life. Proper care of the media can extend
that life. I currently don't have the answers for storage of all forms
of media, and am unlikely to ever have them. It behooves the reader to
determine what the optimal storage environment is for their media. As
most media is designed for use in a typical office environment, it isn’t
a great leap to presume that room temperature with approximately 40-60\%
humidity is likely near optimal. This is in fact the case for the media
that I’ve researched (tapes). All storage media should also be kept in
the dark (light is energy, can be converted to heat, which causes
diffusion) and in as clean an environment as possible. Even for media
that can be cleaned, dirt can not cause more good than harm. It can
eventually end up in places it shouldn’t me.

Further, magnetic media by its very nature should not be exposed to
magnetic fields. Fortunately magnetic fields drop rapidly away from the
source so keeping the media at least 2” away from a magnetic field in a
typical office should be sufficient. It is typically easy to guarantee
multiple feet, making any nearby field irrelevant.

Often the technology that writes to and reads from the media is more
finite than the storage lie of the media. The 5$\frac{1}{4}$”
and 3” floppy combined had a technology life of approximately 20 years
total before being replaced by newer formats with higher storage
densities. Many older scientists and engineers have boxes of programs on
punch cards that can no longer be read.\DUfootnotemark{id16}{id19}{2} As new technology becomes
available, it is prudent to occasionally, say every 10 years, create an
additional archive of the data in the new format (transcription). It is
prudent to understand if that format is useful for archiving or not
(long life, inexpensive, sufficiently popular that there will be readers
in the future, sufficiently robust to errors). The original archive
should not be discarded. The benefit of discarding it (minimal) is
likely heavily outweighed by the cost necessary to verify that the new
media is a perfect duplicate, and will also survive as long.

Further, as archived data is important, so is its physical protection.
Buildings burn down, are flooded, robbed, boxes are lost and accidents
happen. It is prudent to have a second archive in another physical
location. Rarely will this ever be needed. However, if the data is
important, it is certainly worth the redundancy. Most people can use
their home and office for such locations. I've had to use secondary
archives due to failed magnetic media in the past. When one loses data
it can be very trying.

While this document is intended to give some insight as to how to
archive data, it is not intended to cover backups, which are somewhat
like short-term archives that can be replaced. It is left to the reader
to investigate the importance of backing up. As a quick summary, weekly
backups to an external drive, or better yet, an online backup system
(automatically off-site), is prudent. Depending on the rate of your
work, quarterly or so archives are also prudent. Just as in archiving,
the average user tends to lack sufficient paranoia.


\section{Summary%
  \label{summary}%
}

It's your data, and your graphics. I haven't found anyone yet who has
disagreed with this advice,\DUfootnotemark{id17}{id20}{3} and most have immediately told stories of
countless hours wasted compensating for a lack of proper archiving
practice. Data must be stored and documented in a meaningful way that is
absolutely clear. Regarding graphics, use EPS or PDF for all graphics
that are not a pure photographs. Converting a photograph to a JPEG or
TIFF is a simple task on a Mac using the Preview application. On other
platforms there is free software available, or Adobe Acrobat~can be
used. This document is focused more on \textquotedbl{}what\textquotedbl{} than \textquotedbl{}how.\textquotedbl{} Modern
internet search engines are more than adequate for finding out how to
perform these tasks, and they will evolve as tools evolve. While this document is far from exhaustive, it will
set you on the right path to protecting your hard work for future use.


\section{Bibliography%
  \label{bibliography}%
}
%
\DUfootnotetext{id18}{id3}{1}{%
Using GNU software protects you from much of this by being free for
use in any environment.
}
%
\DUfootnotetext{id19}{id16}{2}{%
Ironically, it is only recently that technology is developing
electronic storage media with lives longer than cards!
}
%
\DUfootnotetext{id20}{id17}{3}{%
If you do, please help me improve it!
}

\end{document}
